\documentclass[10pt,a4paper]{amsart}
\usepackage[margin=19mm]{geometry}
\usepackage{amsmath,amssymb,mathtools}
\usepackage[scr=rsfs,cal=euler]{mathalfa}
\usepackage{xfrac}
\usepackage[unicode,hidelinks,bookmarks=false]{hyperref}
\newcommand{\ie}{i.e.\ }
\newcommand{\cf}{cf.\ }
\newcommand{\resp}{respectively\ }
\newcommand{\Subsection}{\subsection}
\providecommand{\nobreakdash}{\nobreak\hbox{-}}
\setlength{\emergencystretch}{2em}
\multlinegap0pt
\usepackage{amssymb}
\usepackage[shortlabels]{enumitem}
\newtheorem{lem}{Lemma}
\def \e {{\bf e}}
\newcommand{\ob}[1]{\left(#1\right)}
\def \u {{\bf u}}
\def \v {{\bf v}}
\title{\it Pretty good plus state transfer in cycles}
\author{Sarojini Mohapatra, Hiranmoy Pal}
\date{Source: arXiv:2603.17595v1 (CC BY 4.0)}
\begin{document}
\maketitle

\noindent\emph{Source and licence.} \it Pretty good plus state transfer in cycles. Version arXiv:2603.17595v1 (CC BY 4.0), distributed by arXiv under the Creative Commons Attribution 4.0 International licence, http://creativecommons.org/licenses/by/4.0/, as recorded in the arXiv licence metadata for this exact version and shown on the abstract page https://arxiv.org/abs/2603.17595v1. Full licence notice: \texttt{licenses/arxiv-2603.17595-CC-BY-4.0.txt}.\par\smallskip

\noindent\textit{Editorial scope.} Counted unit: the article's lem 5l10 (source0.tex lines 852--854). The article's complete original proof of it is delivered with it (source0.tex lines 856--872, one \texttt{proof} environment of 17 lines). No mathematics is written, repaired or re-derived: the statement and the proof are the article's own lines, in the article's own notation, and no retained line is substituted or re-flowed.  Retained uncounted as the source-local context the proof needs: lines 153--155.  Nothing was omitted from any retained span, so the package carries no omission record.  No retained line contains a citation, so the wrapper carries no bibliography and no bibliography entry.  No retained line contains a figure environment, an image-inclusion command or drawing code, so no image is delivered and the package declares no asset dependency.  Editorial formatting only, in addition: an AMS \texttt{amsart} preamble that restates the article's own declarations and macros used by the retained lines, namely the environments \texttt{lem} on separate counters, together with the standard packages those lines need.  The retained environments print in this document as Lemma 1 in order of appearance, whereas the article prints the same items with the numbering of its own full document.  The complete native source of the article as distributed in the arXiv e-print for this exact version is bundled unchanged as \texttt{sources/196630/source0.tex}.
\begin{equation}\label{5eq2}
       \lim_{k \to \infty}U_{M(G)}\ob{t_k}\u=\gamma \v.
   \end{equation}

\begin{lem}\label{5l10}
    Let $n$ be even and $a$ and $b$ be antipodal vertices in the circulant graph $\text{Cay}(\mathbb{Z}_n, S).$ If there is pretty good plus state transfer from $\frac{1}{\sqrt{2}}\ob{\e_a+\e_b},$ then $n$ is divisible by $4$ and pretty good plus state transfer occurs between $\frac{1}{\sqrt{2}}\ob{\e_a+\e_b}$ and $\frac{1}{\sqrt{2}}\ob{\e_{a+\frac{n}{4}}+\e_{b+\frac{n}{4}}}.$  
\end{lem}

\begin{proof}
    Without loss of generality, let $a=0$ and $b=\frac{n}{2}.$ Suppose $\text{Cay}(\mathbb{Z}_n, S)$ admits plus PGST between $\frac{1}{\sqrt{2}}\ob{\e_0+\e_{\frac{n}{2}}}$ and $\frac{1}{\sqrt{2}}\ob{\e_c+\e_d}.$Then by \eqref{5eq2}, we have 
\begin{equation}\label{5e13}
       \lim_{k \to \infty}U\ob{t_k}\ob{\e_0+\e_{\frac{n}{2}}}=\gamma\ob{\e_c+\e_d}.
   \end{equation}
   % Let $P$ be a permutation matrix corresponding to an automorphism of $\text{Cay}(\mathbb{Z}_n, S)$ satisfying $P\e_j=\e_{j+\frac{n}{2}}.$ Multiplying  
    Let $P$ be the permutation matrix corresponding to the automorphism of $\text{Cay}(\mathbb{Z}_n, S)$ satisfying $P\e_j=\e_{j+\frac{n}{2}}.$ Since $P(\e_0+\e_{\frac{n}{2}})=\e_0+\e_{\frac{n}{2}},$ and the limit in \eqref{5e13} is unique, we have
$P(\e_c+\e_d)=\e_c+\e_d.$
Therefore, $c$ and $d$ are antipodal vertices. Let $P'$ be the permutation matrix corresponding to the automorphism of $\text{Cay}(\mathbb{Z}_n, S)$ such that $P'\e_j=\e_{n-j} $ that fixes only $0$ and $\frac{n}{2}.$ Multiplying $P'$ on both sides of \eqref{5e13}, gives 
\begin{equation}\label{5e15}
P'(\e_c+\e_d)=\e_c+\e_d.
\end{equation}
The only antipodal vertices satisfying $\eqref{5e15},$ are $c=\frac{n}{4}$ and $d=\frac{3n}{4}.$ 
% Since 
% $a$ and $b$ are arbitrary antipodal vertices and $\text{Cay}(\mathbb{Z}_n, S)$
% is vertex transitive, the result follows.
\end{proof}

\end{document}
